\originalchapter{非惯性系与相对论接口}\label{ch:frames}
\originalsection{平移与转动参考系}
以加速度 $\vct A$ 平移的系中, 质点相对加速度满足 $m\vct a'=\vct F-m\vct A$. 惯性力 $-m\vct A$ 不是另一物体施加的相互作用, 而是把运动方程改写到该系所需的项. 匀加速车厢中的摆沿有效重力 $\vct g-\vct A$ 平衡.

对绕某原点以 $\vct\Omega(t)$ 转动的坐标系, 基矢的惯性导数为 $\vct\Omega\times\uv e_i$. 任意向量满足
\begin{equation}
(\mathrm d\vct Q/\mathrm dt)_{\rm inertial}
=(\mathrm d\vct Q/\mathrm dt)_{\rm rotating}+\vct\Omega\times\vct Q.
\end{equation}
对位置求两次导数并加入原点加速度, 得
\begin{align}
m\vct a'={}&\vct F-m\vct A-2m\vct\Omega\times\vct v'
-m\vct\Omega\times(\vct\Omega\times\vct r')
-m\dot{\vct\Omega}\times\vct r'.\label{eq:rotframe}
\end{align}
依次为平移惯性力、科里奥利力、离心力、欧拉力. 科里奥利力与相对速度垂直, 在旋转系中瞬时不做功; 离心力可在恒定 $\Omega$ 时用势能 $-m|\vct\Omega\times\vct r'|^2/2$ 表达. 转速变化时漏掉欧拉力会使收臂转动等问题不守恒.

\originalsection{地球及旋转台}
地球角速度 $\Omega_E$, 纬度 $\lambda$. 离心加速度沿距自转轴的径向, 大小 $\Omega_E^2R_E\cos\lambda$. 其地表竖直向外分量为 $\Omega_E^2R_E\cos^2\lambda$, 故有效重力的竖直分量近似 $g-\Omega_E^2R_E\cos^2\lambda$. 纬度不同和实际地球非球形都会修正 $g$, 不能只归因于距离球心不同.

局部东、北、上为右手基, 地球转动向量为 $\Omega_E(\cos\lambda\uv n+\sin\lambda\uv z)$. 水平向北速度 $v\uv n$ 的科里奥利加速度向东, 大小 $2\Omega_Ev\sin\lambda$. 因而北半球水平运动偏向运动方向右侧, 南半球左侧; 自由下落也会受到水平偏转, 必须代入完整局部速度求出, 不能只记一个口号.

惯性系与转台轨迹由旋转矩阵直接联系: 若台角为 $\Omega t$, 则 $x'=x\cos\Omega t+y\sin\Omega t$, $y'=-x\sin\Omega t+y\cos\Omega t$. 惯性系的直线轨迹可在台上看成曲线, 虚拟力不会改变这组几何变换. 转台中心的小角摆, 在水平向量 $\vct r'$ 上有
\begin{equation}
\ddot{\vct r'}+2\Omega\uv z\times\dot{\vct r'}+(\omega_0^2-\Omega^2)\vct r'=0.
\end{equation}
复坐标 $w=x'+iy'$ 满足 $\ddot w+2i\Omega\dot w+(\omega_0^2-\Omega^2)w=0$, 令 $w=e^{-i\Omega t}u$ 化为 $\ddot u+\omega_0^2u=0$. 这说明惯性系摆动平面固定, 相对转台以反向角速度转动. 地球小角摆近似只留竖直自转分量, 得福科摆进动 $-\Omega_E\sin\lambda$, 并需忽略更小离心、垂直耦合与阻尼效应.

\originalsection{狭义相对论的基本限制}
伽利略变换 $x'=x-Vt,t'=t$ 与低速牛顿理论相容, 但不保持所有惯性观察者测得的光速相同. 狭义相对论以惯性系物理规律相同、真空光速相同为基本事实. 沿 $x$ 轴相对速度 $V$ 时洛伦兹变换为
\begin{equation}
x'=\Gamma(x-Vt),\quad t'=\Gamma(t-Vx/c^2),\quad
\Gamma=(1-V^2/c^2)^{-1/2}.
\end{equation}
可从线性变换、相对运动对称性与光线 $x=\pm ct$ 保持的条件确定系数. 事件间隔 $c^2\Delta t^2-\Delta x^2-\Delta y^2-\Delta z^2$ 保持不变. 一只钟在其静止系的两次事件 $\Delta x'=0$, 则运动系时间为 $\Delta t=\Gamma\Delta\tau$, $\Delta\tau$ 为固有时. 同时性的判定改变, 所以测运动杆长度时必须在同一观察系同时取两端坐标, 得 $L=L_0/\Gamma$.

速度变换来自位置与时间微分之比: $u'_x=(u_x-V)/(1-u_xV/c^2)$, 从而不会用简单相加把光速变成 $c+V$. 相对论动量与能量为 $\vct p=\gamma m\vct v$, $E=\gamma mc^2$, 满足 $E^2=p^2c^2+m^2c^4$. $v\ll c$ 时 $E=mc^2+mv^2/2+O(v^4/c^2)$, 恢复牛顿动能. 这些公式在此作为与后续理论物理的接口; 完整协变动力学接相对论或电动力学册, 不把本节当成其全部重编. 不能从某一黑洞质量或半径估算单独得出粒子实验的风险结论.
